\documentclass[10pt,a4paper]{amsart}
\usepackage[margin=19mm]{geometry}
\usepackage{amsmath,amssymb,mathtools}
\usepackage[scr=rsfs,cal=euler]{mathalfa}
\usepackage{xfrac}
\usepackage[unicode,hidelinks,bookmarks=false]{hyperref}
\newcommand{\ie}{i.e.\ }
\newcommand{\cf}{cf.\ }
\newcommand{\resp}{respectively\ }
\newcommand{\Subsection}{\subsection}
\providecommand{\nobreakdash}{\nobreak\hbox{-}}
\setlength{\emergencystretch}{2em}
\multlinegap0pt
\usepackage{xcolor}
\usepackage{amssymb}
\usepackage{tikz}
\newtheorem{lemma}{Lemma}
\newtheorem{theorem}{Theorem}
\def\ni{\noindent}
\title{Mutual-visibility coloring of graphs}
\author{Saneesh Babu, Gabriele Di Stefano, Aparna Lakshmanan S}
\date{Source: arXiv:2512.12251v2 (CC BY 4.0)}
\begin{document}
\maketitle

\noindent\emph{Source and licence.} Mutual-visibility coloring of graphs. Version arXiv:2512.12251v2 (CC BY 4.0), distributed by arXiv under the Creative Commons Attribution 4.0 International licence, http://creativecommons.org/licenses/by/4.0/, as recorded in the arXiv licence metadata for this exact version and shown on the abstract page https://arxiv.org/abs/2512.12251v2. Full licence notice: \texttt{licenses/arxiv-2512.12251-CC-BY-4.0.txt}.\par\smallskip

\noindent\textit{Editorial scope.} Counted unit: the article's theorem binary (source0.tex lines 272--278). The article's complete original proof of it is delivered with it (source0.tex lines 280--308, one \texttt{proof} environment of 29 lines). No mathematics is written, repaired or re-derived: the statement and the proof are the article's own lines, in the article's own notation, and no retained line is substituted or re-flowed.  Retained uncounted as the source-local context the proof needs: lines 257--259.  Nothing was omitted from any retained span, so the package carries no omission record.  No retained line contains a citation, so the wrapper carries no bibliography and no bibliography entry.  No retained line contains a figure environment, an image-inclusion command or drawing code, so no image is delivered and the package declares no asset dependency.  Editorial formatting only, in addition: an AMS \texttt{amsart} preamble that restates the article's own declarations and macros used by the retained lines, namely the environments \texttt{lemma}, \texttt{theorem} on separate counters, together with the standard packages those lines need.  The retained environments print in this document as Lemma 1, Theorem 1 in order of appearance, whereas the article prints the same items with the numbering of its own full document.  The complete native source of the article as distributed in the arXiv e-print for this exact version is bundled unchanged as \texttt{sources/190927/source0.tex}.
\begin{lemma} \label{lem}
If $S$ is a mutual-visibility set of $GT(r)$, then $|{S \cap V(C_{i,j}) }| \leq 3 $.
\end{lemma}

\begin{theorem} \label{binary}
For $r \geq 2$,
	$$\chi_\mu(GT(r))=\begin{cases}
        2(r-i) + 3, & 2^{i-1}+i-2 \leq r \leq 2^{i-1}+2^{i-2}+i-3\\
	    2(r-i) + 2, & 2^{i-1}+2^{i-2}+i-2 \leq r \leq 2^i + i - 2       
	\end{cases}$$
\end{theorem} 

 \begin{proof}
We begin the proof by establishing a coloring of $GT(r)$ using the number of colors specified in the theorem.\par

\ni{\bf Case 1: $2^{i-1}+2^{i-2}+i-2 \leq r \leq 2^i + i - 2$}.\\
In this case, assign color $1$ to all the quasi-leaves, color $2k$ to all the vertices of depth $r-k$ in $V_r^{(1)}$ and to the $k^{th}$ vertex in the sequence $\{v'_{1,1}, v'_{2,1},v'_{2,2}, v'_{3,1},$ $v'_{3,2}, v'_{3,3}, v'_{3,4},\ldots\}$, for $k = 1, 2, \ldots, r-i$. Similarly, assign color $2k+1$ to all vertices of depth $r-k$ in $V_r^{(2)}$ and to the $k^{th}$ vertex in the sequence $\{v_{1,1}, v_{2,2},v_{2,1}, v_{3,4}, v_{3,3}, v_{3,2}, v_{3,1},\ldots\}$, for $k = 1, 2, \ldots, r-i$. Here all vertices of depth $0,1,\ldots,i-2$ and $r-1,r-2,\ldots,r-(r-i)=i$ are completely colored and vertices of depth $i-1$ are partially colored. Among the vertices of depth $i-1$ in $V_r^{(2)}$, $v'_{i,1},v'_{i,2},\ldots,v'_{i,j}$ will be already colored, where $j=r-i-(2^{i-1}-1) \geq 2^{i-2}-1$, (since, $r \geq 2^{i-1}+2^{i-2}+i-2$). Similarly, in $V_r^{(1)}$, $v_{i,a},v_{i,a-1},\ldots,v_{i,a-j-1}$ will be already colored, where $a=2^{i-1}$ and $j \geq 2^{i-2} -  1$. When $r \geq 2^{i-1}+2^{i-2}+i - 1$, we have $j \geq 2^{i-2}$ so that there are fewer than $2^{i-2}$ uncolored vertices of depth $i-1$. In this case, color all the remaining uncolored vertices with color $2(r-i)+2$. When $r = 2^{i-1}+2^{i-2}+i - 2$, color $v_{i-1,j+1}$ with color 1 and all the remaining vertices with color $2(r-i)+2$ to obtain a mutual-visibility coloring of $GT(r)$ with $2(r-i)+2$ colors.\\
\ni{\bf Case 2: $2^{i-1}+i-2 \leq r \leq 2^{i-1}+2^{i-2}+i-3$}.\\
\ni{\bf Sub Case 1: $r=2^{i-1}+i-2$}.\\
In this case, we proceed with the same coloring pattern given in Case 1 up to $k=r-i+1$ so that the entire vertices of $GT(r)$ are colored using  $2(r-i)+3$ colors.\\
\ni{\bf Sub case 2: $2^{i-1}+i-1 \leq r \leq 2^{i-1}+2^{i-2}+i-3$}.\\
The coloring procedure is the same as Case 1, where the only difference is that more than half of the vertices of depth $i-1$ (all vertices when $r=2^{i-1}+i-1$) will be uncolored. We can use one color each for these uncolored vertices in $V_r^{(1)}$ and $V_r^{(2)}$ so that the resulting coloring is a mutual-visibility coloring of $GT(r)$ by Lemma \ref{lem}. 

To prove the converse, first we prove that $2(r-i)+1$ colors are not enough for a mutual-visibility coloring of $GT(r)$, where $i$ is the unique integer such that $2^{i-1} + i-2 \leq r \leq 2^i+i-2$. Assume the contrary that there exists a mutual-visibility coloring with $2(r-i)+1$ colors.\par

Consider the cycle $C_{1,2^r}$. By Lemma \ref{lem}, a color can be used at most 3 times to color the vertices of this cycle. Since we have only $2(r-i)+1$ colors, and there are $4r$ vertices in $C_{1,2^r}$, at least $4r - 2[2(r-i)+1] = 4i-2$ colors must be used at least three times. Let $\chi_0$ be the collection of colors which are used three times to color the vertices of $C_{1,2^r}$ so that $|\chi_0|\geq4i-2$. Now, consider the cycle $C_{2^{r-1},2^{r-1}+1}$, which is also a cycle of length $4r$. Note that, $|V(C_{1,2^r}) \cap V(C_{2^{r-1},2^{r-1}+1})| = 6$ and there are $4r-6$ vertices in $V(C_{2^{r-1},2^{r-1}+1}) \setminus V(C_{1,2^r}) $. Similar to the discussion above, at least $4i - 8$ colors  must be used three times to color these vertices. Let $\chi_1$  be the collection of colors which are used three times to color the vertices in $V(C_{2^{r-1},2^{r-1}+1}) \setminus V(C_{1,2^r})$ so that $|\chi_1|\geq4i - 8$. We claim that $\chi_0  \cap \chi_1  = \phi$. On the contrary, if $c \in \chi_0  \cap \chi_1$ then $c$ must be used at least two times in the geodesic path $v_{1,1}-v_1-v'_{1,1}$ or in the geodesic path $v_{1,1}-v_{2^r}-v'_{1,1}$. Similarly, $c$ must be used at least two times in the geodesic $v_{3,2}-v_{2^{r-1}}-v'_{3,2}$ or $v_{3,3}-v_{2^{r-1}+1}-v'_{3,3}$. But then, in one of the cycles $C_{1,2^{r-1}}, C_{1,2^{r-1}+1}, C_{2^{r-1},2^r}$ or $C_{2^{r-1}+1,2^r}$, $c$ will be used four times which contradicts Lemma \ref{lem}. Therefore,$\chi_0  \cap \chi_1  = \phi$.


Now, consider the cycles 
%$C_{2^{r-2},2^{r-2}+1}$ and $C_{2^{r-1}+2^{r-2},2^{r-1}+2^{r-2}+1}$
$C_{t2^{r-2},t2^{r-2}+1}$, for $t=1,3$. These two cycles have $4r-12$ vertices each which are not part of both the previous cycles considered. Similar to the case above we can find a set of colors $\chi_2$ which are used three times either in $V(C_{2^{r-2},2^{r-2}+1}) \setminus \{ V (C_{1,2^r})\cup V( C_{2^{r-1},2^{r-1}+1})\}$ and $V(C_{2^{r-1}+2^{r-2},2^{r-1}+2^{r-2}+1}) \setminus \{V(C_{1,2^r}) \cup V(C_{2^{r-1},2^{r-1}+1})\}$. Also, $ \chi_2 \cap \chi_1 = \chi_2 \cap \chi_0 = \phi$ and $|\chi_2| \geq 2(4i - 12)$. Now, we consider the four cycles 
%$C_{t+2^{r-3},t+2^{r-3}+1}$, where $t = 0, 2^{r-2}, 2^{r-1}, 2^{r-1}+2^{r-2}$ 
$C_{t2^{r-3},t2^{r-3}+1}$, for $t=1,3,5,7$, and proceed as above to get $\chi_3$ of cardinality at least $4(4i - 16)$ which do not intersect pairwise with any of $\chi_0,\chi_1$ and $\chi_2$. Continue like this to get a collection of distinct colors $\chi_1, \chi_2,\ldots,\chi_{i-1}$, where each $\chi_j$ is of cardinality at least $2^{i-1}(4i - 4j - 4)$, $j \in [i-1]$ and $\chi_0$ of cardinality at least $4i-2$. Note that the disjoint sets $\chi_j$, $j \in [i-1]$ have cardinality at least $2^{i-1}(4i - 4j - 4)$ which is an arithmetic-geometric progression with common difference -4 and common ratio 2. Therefore, $|\cup_{j=1}^{i-1}\chi_j |\geq 2^{i+1}-4i$ and together with the $4i-2$ more distinct colors in $\chi_0$, at least $2^{i+1}-2$ colors are used three times to color $GT(r)$. But, we have only $2(r-i)+1$ colors and $r\leq 2^i+i-2$ so that, $2(r-i)+1 \leq 2^{i+1}-3$, a contradiction. 

%\textcolor{blue}{I do not understand the last two sentences from ``Therefore, ....''. The distinct colors in $C_0$ are at least $4i-2$ (not $4i$ as stated) for a total of at least $2^{i+1}-6$ (not $2^{i+1}-2$) colors. This is not in contradiction with the number of colors that is at most  $2^{i+1}-3$.  }\par
%\textcolor{green}{The error in computation is ratified.}

Now, if we do the same procedure with $2(r-i)+2$ colors, then $|\chi_0| \geq 4i - 4$ and $|\chi_j| \geq 2^{i-1}(4i - 4j - 6)$ so that $|\cup_{j=0}^{i-2}\chi_j| \geq 2^{i}+2^{i-1}-2$. When $r \leq 2^{i-1}+2^{i-2} + i - 3$, the number of colors $2(r-i)+2 \leq 2^{i}+2^{i-1}-4$, a contradiction. \par

Hence, the theorem.
\end{proof}

\end{document}
